\section{Prime moduli and arbitrary integer weights}
\begin{lemma}\label{lem:primorial}
If $d,e\geq2$ and $n=de$, every prime $p\leq n$ divides $Q(n,d)$.
\end{lemma}
\begin{proof}
Counting multiples of $p,p^2,\ldots$ in each factorial gives
\[
 v_p(Q(n,d))=\sum_{j\geq1}\left(\left\lfloor\frac{de}{p^j}\right\rfloor
 -\left\lfloor\frac d{p^j}\right\rfloor
 -\left\lfloor\frac e{p^j}\right\rfloor\right).
\]
Every summand is nonnegative since $de\geq d+e$.
Interchange $d,e$ so $d\leq e$. Suppose $p\geq3$.
If $e<p$, the $j=1$ summand is positive because $de\geq p$.
If $e\geq p$, it is at least
$(d-1)\lfloor e/p\rfloor-\lfloor d/p\rfloor
\geq d-1-\lfloor d/p\rfloor\geq1$.
For $p=2$ this inequality works when $d\geq3$.
When $d=2,e\geq3$, the $j=1$ summand is
$e-1-\lfloor e/2\rfloor\geq1$; the remaining case has $Q(4,2)=6$.
\end{proof}

\begin{corollary}\label{cor:prime}
For every prime $p$, integer $u,v$, and $n\geq p$,
\[
 \A(n)\equiv uv^n+vu^n\pmod p,\qquad
 \A(n+p-1)\equiv\A(n)\pmod p.
\]
Thus A258899 has period $10$ from $n=11$ modulo $11$,
and A258903 has period $6$ from $n=7$ modulo $7$.
\end{corollary}
\begin{proof}
Lemma~\ref{lem:primorial} removes every proper-divisor term.
For a base not divisible by $p$, the unit-product argument in
Lemma~\ref{lem:powers} gives period $p-1$; a base divisible by $p$
has zero powers for all positive exponents. These are sufficient onsets.
\end{proof}

\begin{corollary}\label{cor:triangle}
For any integer-valued function $W(d,e)$,
$\sum_{d\mid n}W(d,n/d)Q(n,d)$ is congruent to
$W(1,n)+W(n,1)$ modulo $m$ if $n\geq2m+4$, or modulo a prime
$p$ if $n\geq p$. In A386877 \cite{a386877},
all entries strictly between the endpoint $1$'s vanish modulo
the respective modulus in these ranges; the row sum is $c(n)$ for $n\geq1$.
The sign of $T(n,k)$ is the divisibility indicator A113704 for $1\leq k\leq n$.
\end{corollary}
\begin{proof}
Apply Lemma~\ref{lem:factorial} or Lemma~\ref{lem:primorial} termwise.
Nondivisor entries and column zero are zero for $n>0$;
the other entries are the positive integers $Q(n,d)$ by the triangle's definition,
which also proves the sign assertion.
\end{proof}

\section{Direct cross-references and a history observation}
\begin{corollary}\label{cor:interpretations}
The number of matrices with entries $1,\ldots,n$, each used once,
up to independent row and column permutations is $c(n)$.
The same count enumerates the labeled rectangular-band operations
$(i,j)(k,l)=(i,l)$. Also $c(n)=2$ exactly when $n$ is prime.
\end{corollary}
\begin{proof}
For each shape $d\times e$, $de=n$, the action of $S_d\times S_e$
on the $n!$ distinct-entry arrays is free, giving $Q(n,d)$ orbits.
A grid gives the stated operation. Conversely its row and column classes
are recovered from equality of left and right multiplication maps,
so two grids give the same operation exactly when they differ by row
and column permutations. Sum over the shapes.
For prime $n$ only the two endpoint terms occur; for composite $n>1$
a positive proper-divisor term adds to $2$, and $c(1)=1$.
The rectangular-band enumeration and primality criterion already appear in
Devillet and Kiss \cite[Proposition 4.13 and Corollary 4.15]{dk},
whose preprint dates from June 2018.
\end{proof}

\begin{corollary}\label{cor:quet}
For A057625 \cite{a057625}, let $L(n)=\sum_{d\mid n}n!/d!$, $n\geq1$.
Then $L(n)\equiv1\pmod m$ for every $m\geq1$, $n\geq2m$, and
$\sum_{n\geq1}L(n)x^n/n!=\sum_{d\geq1}x^d/(d!(1-x^d))$.
\end{corollary}
\begin{proof}
The term $d=n$ is $1$. For a proper divisor $d\leq n/2$,
$n-d\geq m$, and
$n!/d!=m!(n-d-m)!\binom{n}{d,m,n-d-m}$ is divisible by $m$.
Expand $x^d/(1-x^d)=\sum_{e\geq1}x^{de}$ for the generating function.
\end{proof}

\begin{corollary}\label{cor:other}
Let $P(n),B(n),F(n)$ denote A038041, A061095 and A236696,
respectively \cite{other}; their offset is one. Then
\begin{align*}
 \sum_{n\geq1}P(n)\frac{x^n}{n!}
   &=\sum_{k\geq1}\left(\exp(x^k/k!)-1\right),\\
 \sum_{n\geq1}B(n)\frac{x^n}{n!}
   &=\sum_{k\geq1}\frac{x^k}{k!-x^k},\\
 \sum_{n\geq1}F(n)\frac{x^n}{n!}
   &=\sum_{k\geq1}\left(\exp(k^{k-1}x^k/k!)-1\right).
\end{align*}
Also $P(n)=2$ exactly at prime $n$, and
$0\leq B(n)/n!-1\leq n2^{-n/2}$.
If $H(n)$ is A320444 \cite{a320444}, its prime specialization is
$H(p+1)=1+(p+1)^{p-1}$ for every prime $p$.
\end{corollary}
\begin{proof}
The defining sums are
\begin{align*}
 P(n)&=\sum_{ke=n}\frac{n!}{e!(k!)^e},&
 B(n)&=\sum_{ke=n}\frac{n!}{(k!)^e},\\
 F(n)&=\sum_{ke=n}\frac{n!}{e!}\left(\frac{k^{k-1}}{k!}\right)^e.
\end{align*}
Expand the exponentials and geometric series; every positive-degree
coefficient is a finite sum. The displayed A236696 generating function
omits the subtractions of $1$; its constant coefficient is an infinite
sum of $1$'s. The formula above is the well-defined offset-one version.
For $P$, the endpoint partitions contribute $2$ when $n>1$;
a proper factor supplies additional partitions at composite $n$, and $P(1)=1$.
For $B$, the $k=1$ term is $n!$. If $k\geq2$, then
$k!\geq2^{k-1}\geq2^{k/2}$, so $(k!)^{n/k}\geq2^{n/2}$;
there are at most $n$ other terms, proving the bound and $B(n)\sim n!$.
Finally
$H(n+1)=\sum_{d\mid n}Q(n,d)((n+1)/d)^{n/d-1}$
has only the divisors $1,p$ when $n=p$, contributing
$(p+1)^{p-1}$ and $1$. Thus the printed prime formula in A320444
needs the shift to $H(p+1)$: $H(3)=4$, whereas $17=1+4^2=H(4)$.
These sums have different denominators; no congruence for $P,B,F,H$
is inferred from Theorem~\ref{thm:main}.
\end{proof}

\begin{corollary}\label{cor:three}
Three consecutive residues $2$ in A121860 need not begin a constant tail.
Modulo $89$, $(c(41),c(42),c(43),c(44))\equiv(2,2,2,76)$.
\end{corollary}
\begin{proof}
The numbers $41,43$ are prime. Pairing complementary proper divisors gives
$c(42)=2+2(Q(42,2)+Q(42,3)+Q(42,6))$ and
$c(44)=2+2(Q(44,2)+Q(44,4))$.
The quotient residues modulo $89$ are $(82,71,25)$ and $(21,16)$,
respectively, giving the asserted residues by exact arithmetic.
This corrects David Pasino's observation in the A121860 history,
15 November 2025, discussion following revision 36 \cite{history}.
\end{proof}

\begingroup\small
\begin{longtable}{@{}>{\raggedright\arraybackslash}p{0.25\textwidth}>{\raggedright\arraybackslash}p{0.39\textwidth}>{\raggedright\arraybackslash}p{0.29\textwidth}@{}}
\toprule
Source item & Statement & Coverage \\
\midrule
\endhead
A121860, COMMENTS (2019, 2025) & Matrix and rectangular-band counts & Proved; prior band proof cited: Corollary~\ref{cor:interpretations} \\
A121860 history, 15 Nov 2025 & Three residues $2$ predict onset & Refuted-and-corrected: Corollary~\ref{cor:three} \\
A386877 & Row sum and vanishing interior entries & Proved: Corollary~\ref{cor:triangle} \\
A057625 & Constant tail $1$; generating function & Proved: Corollary~\ref{cor:quet} \\
A038041 & Prime criterion; generating function & Proved: Corollary~\ref{cor:other} \\
A061095 & Generating function; $B(n)\sim n!$ & Proved: Corollary~\ref{cor:other} \\
A236696, FORMULA & Generating function's constant coefficient & Refuted-and-corrected: Corollary~\ref{cor:other} \\
A320444, FORMULA & Prime-specialization index & Refuted-and-corrected: Corollary~\ref{cor:other} \\
A100195 \cite{a100195} & Record Bernoulli denominators & Not covered: different arithmetic question \\
A038041, A061095, A236696, A320444 & Possible modular tails & Not covered: denominators differ \\
\bottomrule
\end{longtable}
\endgroup
